\section{Planteamiento del Problema}
    Implementar y comparar sobre un conjunto de imágenes:
    \begin{enumerate}
        \item Transformación afín general para un imagen en escala de grises.
        \item Interpolación bilineal
    \end{enumerate}

    Así como probar la implementación sobre una matriz de transformación afín que realice, al menos, una operación de: 
    \begin{enumerate}
        \item Rotación, con un ángulo que no sea múltiplo de 45°.
        \item Traslación en ambos ejes.
        \item Reflexión en ambos ejes.
        \item Escalamiento en ambas direcciones/ejes
        \item Al menos dos ejemplos de composición de transformaciones afines: rotación, traslación, escalamiento (an)isotrópico, reflexión y/o cizallamiento.
    \end{enumerate}

    Las operaciones indicadas deberán ser implementadas como matrices de transformación afín, utilizando la interpolación bilineal para el relleno de huecos.\\

    Ambas operaciones deberán implementarse manualmente con NumPy, sin utilizar una función de alto nivel que haga el procedimiento.